\section{Future work and the next experimental stage}
\label{sec:future}

The v4.1/v5.1 results identify two priorities: path ranking, especially for
reverse reasoning, and the benchmark's low sensitivity to structural
intervention. The next stage should address these two issues separately to
avoid conflating retrieval changes with changes to the verification engine.

\subsection{Query-aware retrieval and path reranking}

Answer intent and the roles of cause, mediator, and effect should be
identified before the primary path is selected. For reverse reasoning, the
outcome detected in the query should be fixed as the endpoint, after which
the system should perform constrained backward search to find a predecessor.
The path reranker may combine semantic scores with lexical and structural
features, including article number, penalty level, event position on the
path, rule condition/effect, and outcome specificity.

Each retrieval change should be evaluated before verification on the same
candidate pool. The primary metrics should include Rule/Event Recall, Oracle
Path, Top-1 Path, and ranking error. Only after retrieval has been frozen
should the paired structural SCM--path ablation experiment be rerun, thereby
preserving the ability to attribute any difference to the verification
engine.

\subsection{An intervention-sensitive benchmark}

A new structural-counterfactual benchmark should represent each sample as:

\begin{equation}
\left(
q,\,
C_F,\,
\doop(M=m),\,
Y_F,\,
Y_{\doop(M=m)},\,
z
\right),
\label{eq:future-counterfactual-item}
\end{equation}

where $q$ is the query, $C_F$ is the factual context, $M$ is the intervened
mediator, $Y_F$ and $Y_{\doop(M=m)}$ are the factual and counterfactual
outcomes, respectively, and $z$ is the relation label. At a minimum, the
benchmark should balance three labels:

\begin{itemize}
    \item \texttt{NECESSARY}: intervening on the mediator changes the outcome;
    \item \texttt{NON\_NECESSARY}: the outcome is maintained by an active
    alternative mechanism;
    \item \texttt{INDETERMINATE}: the facts, root state, or rules are
    insufficient to support a conclusion.
\end{itemize}

The \texttt{SUFFICIENT}, \texttt{NO\_EFFECT}, and conflict groups may be
added to test additional branches of the engine. The benchmark must contain
pairs for which node deletion and hard intervention yield different results;
otherwise, the paired ablation will remain unable to distinguish the
semantics of the two methods.

Each sample should be reviewed by experts in both the factual and intervened
worlds. Gold annotations should include the rule applied, the rule disabled,
the alternative mechanism, the states of the mediator and outcome, and a
natural-language explanation.

\subsection{Extending rule representation and LegalSCM}

The next model should support conjunctive conditions, negation, exceptions,
rule priority, and temporal validity. A general rule may be represented by
multiple literals in its antecedent rather than by a single positive
condition. When conflicts arise, the system must clearly distinguish
\texttt{FALSE} from \texttt{UNKNOWN} and record which rules produce each
state.

The event schema should also separate the canonical ID, aliases, and
preferred display name. Aliases should be used for retrieval, whereas the
preferred name should be used in answers. This change would both reduce
semantic noise and prevent strings containing \texttt{||} from appearing in
the output.

\subsection{Evaluation of nonlinear structures and citations}

Branch and convergence structures should be evaluated with path-set or
subgraph metrics rather than being forced into a single chain. Planned
metrics include branch outcome coverage, convergence input coverage,
subgraph edge F1, and graph edit distance.

For citations, evaluation should move from article-set matching to
assertion-level grounding. Each assertion in the short answer or explanation
must be mapped to a supporting rule and article. Experts can then assess
three separate criteria: whether the article is relevant, whether it entails
the assertion, and whether the citation is sufficient to support the
conclusion.

\subsection{External and human evaluation}

Future experiments should include baselines run with the same corpus and
protocol, including vanilla RAG, graph RAG without verification, LLM-only,
and reproducible causal/counterfactual RAG methods. An independent test set
authored by experts should be separated from the graph-construction pipeline.

Beyond automated metrics, the study should measure whether the LegalSCM
trace actually helps evaluators:

\begin{itemize}
    \item detect an incorrect path or citation more quickly;
    \item understand the difference between factual and counterfactual worlds;
    \item recognize an alternative mechanism;
    \item calibrate their confidence in the answer.
\end{itemize}

\subsection{Strengthening reproducibility}

Before a conclusive run, the code commit, embedding-model revision,
dependency lockfile, random seed, server configuration, and hashes of all
input/output artifacts should be fixed. Both modes must start from the same
cached retrieval output and record paired differences at the sample level.
Runtime should be measured repeatedly under both cold-cache and warm-cache
conditions.

The hypotheses, primary metrics, exclusion rules, and subgroup analyses
should be specified before the final results are examined. This organization
helps prevent post hoc metric selection and clarifies whether an improvement
comes from retrieval, verification, or answer generation.
